\documentclass[UTF8]{ctexart}
\usepackage{xeCJK}
\usepackage{geometry}
\geometry{left=2cm,right=2cm,top=2.5cm,bottom=2.5cm}
\setlength{\headheight}{12.7pt}

\usepackage{titlesec}
\titleformat{\section}
  {\normalfont\large\bfseries}
  {\thesection}
  {1em}
  {}
\titlespacing*{\section}
  {0pt}
  {3.5ex plus 1ex minus .2ex}
  {2.3ex plus .2ex}

\usepackage{amsmath}
\usepackage{amssymb}
\usepackage{amsthm}
\usepackage{enumitem}
\setlist[enumerate]{itemsep=0pt,topsep=0pt,parsep=0pt,leftmargin=0pt,itemindent=2.5em}
\setlist[itemize]{itemsep=0pt,topsep=0pt,parsep=0pt,leftmargin=0pt,itemindent=2.5em}
\usepackage{fancyhdr}
\pagestyle{fancy}
\fancyhf{}
\usepackage{graphicx}
\usepackage{hyperref}
\usepackage{indentfirst}
\usepackage{lmodern}

\begin{document}
\setlength{\abovedisplayskip}{1pt}
\setlength{\belowdisplayskip}{1pt}

\section{势关系}

\hypertarget{sec:定义1.1}{}
\noindent\textbf{定义1.1}

给定集合$a,b$, 一般使用下面的记号:
\begin{enumerate}
    \item 如果存在从$a$到$b$的\href{03 映射#nameddest=sec:定义2.2}{双射},
    就记$a\approx b$, 称$a,b$\textbf{等势},
    \item 如果存在从$a$到$b$的\href{03 映射#nameddest=sec:定义2.1}{单射},
    就记$a\preccurlyeq b$,
    \item 如果$a\preccurlyeq b$且$a\not\approx b$,
    就记$a\prec b$.
\end{enumerate}

\vspace{15pt}

\hypertarget{sec:定理1.1}{}
\noindent\textbf{定理1.1}

任给集合$a,b$, 如果$b$非空, 那么$a\succcurlyeq b$当且仅当
存在从$a$到$b$的\href{03 映射#nameddest=sec:定义2.2}{满射}.

\noindent{\kaishu 证明.}

\noindent{\kaishu 必要性.}

取单射$f:b\to a$. 再取$x_0\in b$, 然后对任意的$y\in a$定义
\vspace{3pt}
\[g(y)=\left\{\begin{aligned}
    &f^{-1}(y),\quad&&y\in\mathrm{Ran}\,f,\\[-3pt]
    &x_0,\quad&&y\notin\mathrm{Ran}\,f,
\end{aligned}\right.\\[3pt]\]
则$g$的值域是$f^{-1}[\,\mathrm{Ran}\,f\,]\cup\{x_0\}=b$,
即$g$是从$a$到$b$的满射.

\noindent{\kaishu 充分性.}

取满射$g:a\to b$. 应用\href{03 映射#nameddest=sec:定理3.2}{单射定理},
存在单射$h$使得$h\subset g$并且
$\mathrm{Ran}\,h=\mathrm{Ran}\,g=b$.

考虑$f=h^{-1}$, 则$f$也是单射, 并且
\[\left\{\begin{aligned}
    &\mathrm{Dom}\,f=\mathrm{Ran}\,h=b,\\[-3pt]
    &\mathrm{Ran}\,f=\mathrm{Dom}\,h\subset\mathrm{Dom}\,g=a,
\end{aligned}\right.\]
即$f$是从$b$到$a$的单射, $a\succcurlyeq b$. \hfill $\qedsymbol$

% \vspace{10pt}

% \hypertarget{sec:定理1.2}{}
% \noindent\textbf{定理1.2}

% 任给集合$a,a',b,b',c$, 假设$a\approx a'$, $b\approx b'$, 则:
% \begin{enumerate}
%     \item $\mathcal{P}(a)\approx\mathcal{P}(a')$,
%     \item 如果$a\cap b=a'\cap b'=\varnothing$, 那么$a\cup b\approx a'\cup b'$,
%     \item $a\times b\approx b\times a$,
%     \item $a\times b\approx a'\times b'$,
%     \item $^ab\approx\,^{a'}b'$,
%     \item 如果$a\cap b=\varnothing$, 那么$^{a\cup b}c\approx\,^ac\times\,\!^bc$,
%     \item $^{a\times b}c\approx\,\!^a(\,^bc)$.
% \end{enumerate}

% \noindent{\kaishu 证明.}

% 这里不给出完整证明, 只把双射$\varphi$构造出来. 设$f:a\to a'$和$g:b\to b'$是双射.

% \begin{enumerate}
%     \item $\varphi:\mathcal{P}(a)\to\mathcal{P}(a'),\,s\mapsto f[s]$.
%     \hspace{10.01em}$\varphi^{-1}:s'\mapsto f^{-1}[s']$.
%     \vspace{6pt}
%     \item $\varphi:a\cup b\to a'\cup b',\,x\mapsto\left\{\begin{aligned}
%     &f(x),&&x\in a,\\&g(x),&&x\in b.\end{aligned}\right.$
%     \hspace{4.96em}$\varphi^{-1}:x'\mapsto\left\{\begin{aligned}
%     &f^{-1}(x'),&&x'\in a',\\&g^{-1}(x'),&&x'\in b'.\end{aligned}\right.$
%     \vspace{6pt}
%     \item $\varphi:a\times b\to b\times a,\,(x,y)\mapsto(y,x)$.
%     \hspace{7.55em}$\varphi^{-1}:(y,x)\mapsto(x,y)$.
%     \item $\varphi:a\times b\to a'\times b',\,(x,y)\mapsto(f(x),g(y))$.
%     \hspace{4.35em}$\varphi^{-1}:(x',y')\mapsto(f^{-1}(x'),g^{-1}(y'))$.
%     \item $\varphi:\,^ab\to\,\!^{a'}b',\,\psi\mapsto g\circ\psi\circ f^{-1}$.
%     \hspace{8.67em}$\varphi^{-1}:\psi'\mapsto g^{-1}\circ\psi'\circ f$.
%     \item $\varphi:\,^{a\cup b}c\to\,^ac\times\,\!^bc,\,
%     \psi\mapsto(\psi\!\restriction_a,\psi\!\restriction_b)$.
%     \hspace{6.36em}$\varphi^{-1}:(\psi_1,\psi_2)\mapsto\psi_1\cup\psi_2$.
%     \item $\varphi:\,^{a\times b}c\to\,\!^a(\,^bc),\,\psi\mapsto
%     (x\mapsto(y\mapsto\psi(x,y)))$.
%     \hspace{2.86em}$\varphi^{-1}:\psi'\mapsto((x,y)\mapsto\psi'(x)(y))$.
%     \hfill $\qedsymbol$
% \end{enumerate}

\vspace{15pt}

\hypertarget{sec:定理1.2}{}
\noindent\textbf{定理1.2 (Cantor定理)}

任给集合$a$, 有$a\prec\mathcal{P}(a)$.

\noindent{\kaishu 证明.}

首先, 对任意的$x\in a$定义
\[
    f(x)=\{x\},
\]
则$f:a\to\mathcal{P}(a)$, 并且显然是单射. 所以$a\preccurlyeq\mathcal{P}(a)$.

其次, 假设$a\approx\mathcal{P}(a)$, 即存在双射$g:a\to\mathcal{P}(a)$. 此时定义
\[
    b=\{x\in a\mid x\notin g(x)\},
\]
则$b\in\mathcal{P}(a)$, 再令$y=g^{-1}(b)$. 依定义,
\[
    y\in b\iff y\in g(y)\iff y\notin b,
\]
矛盾. 所以$a\not\approx\mathcal{P}(a)$.

综上, $a\prec\mathcal{P}(a)$. \hfill $\qedsymbol$





\newpage

\hypertarget{sec:定理1.3}{}
\noindent\textbf{定理1.3 (Cantor-Bernstein定理)}

任给集合$a,b$, 如果$a\preccurlyeq b$且$b\preccurlyeq a$, 那么$a\approx b$.

\noindent{\kaishu 证明.}

取单射$f:a\to b$和$g:b\to a$, 记$h=g\circ f$, $a'=a\setminus\mathrm{Ran}\,g$.

定义$a$的子集族
\vspace{-3pt}
\[
    \mathcal{A}=\{\,U\subset a\mid a'\cup h[U]\subset U\,\},\\[-3pt]
\]
显然$a\in\mathcal{A}$, 即$\mathcal{A}$非空.
考虑$\displaystyle V=\bigcap\mathcal{A}$. 首先
\vspace{3pt}
\[\left\{\begin{aligned}
    \,&\,\,\,\,V\!\!\!\!&&=\,\,\,\,\bigcap_{U\in\mathcal{A}}U\!\!\!\!&&
    \subset\bigcap_{U\in\mathcal{A}}a\!\!\!\!&&=a\,,\\
    \,&\,\,\,\,a'\!\!\!\!&&=\,\,\,\,\bigcap_{U\in\mathcal{A}}a'\!\!\!\!&&
    \subset\bigcap_{U\in\mathcal{A}}U\!\!\!\!&&=V\,,\\
    \,&h[V]\!\!\!\!&&=\bigcap_{U\in\mathcal{A}}h[U]\!\!\!\!&&
    \subset\bigcap_{U\in\mathcal{A}}U\!\!\!\!&&=V\,,
\end{aligned}\right.\\[3pt]\]
综上即得$V\subset a$且$a'\cup h[V]\subset V$, 即$V\in\mathcal{A}$.

\vspace{3pt}
其次由$a'\cup h[V]\subset V$得
\[
    h[a'\cup h[V]]\subset h[V]\quad\implies\quad
    a'\cup h[a'\cup h[V]]\subset a'\cup h[V]\quad\implies\quad
    a'\cup h[V]\in\mathcal{A},
\]
又因为$V$是$\mathcal{A}$的交集, 所以$V\subset a'\cup h[V]$.

\vspace{3pt}
综上$V=a'\cup h[V]$. 据此定义映射
\vspace{-3pt}
\[p:a\to b,\quad x\mapsto\left\{\begin{aligned}
    &f(x),\quad&&x\in V,\\[-3pt]
    &g^{-1}(x),\quad&&x\notin V,
\end{aligned}\right.\\[3pt]\]
当$x\notin V$时有$x\notin a'$, 即$x\in\mathrm{Ran}\,g$, 所以$g^{-1}(x)$是存在的.
下证$p:a\to b$是双射, 从而$a\approx b$.

\vspace{3pt}
假设$p(x)=p(x')$, 如果$x,x'$同时属于$V$或不属于$V$, 显然都有$x=x'$.
如果$x\in V,\,x'\notin V$, 那么
\[
    f(x)=g^{-1}(x')\quad\implies\quad x'=h(x)\in h[V]
    \quad\implies\quad x'\in V,
\]
矛盾. 于是$p$是单射.

\vspace{3pt}
任取$y\in b$, 如果$g(y)\in V$则一定有$g(y)\in h[V]$, 从而存在$x\in V$使得
\[
    g(y)=h(x)\quad\implies\quad y=f(x)\quad\implies\quad y=p(x),\\[3pt]
\]
如果$g(y)\notin V$则$y=g^{-1}(g(y))=p(g(y))$. 于是$p$是满射. \hfill $\qedsymbol$

\vspace{20pt}

最后, 可以完整描述$\approx$的等价性和$\preccurlyeq$的偏序性.

\vspace{10pt}

\hypertarget{sec:定理1.4}{}
\noindent\textbf{定理1.4}

任给集合$a,b,c$ :
\begin{enumerate}
    \item $a\approx a$, $a\preccurlyeq a$,
    \item 如果$a\approx b$, 那么$b\approx a$,
    \item 如果$a\approx b$且$b\approx c$, 那么$a\approx c$,
    \item 如果$a\preccurlyeq b$且$b\preccurlyeq c$, 那么$a\preccurlyeq c$,
    \item 如果$a\preccurlyeq b$且$b\preccurlyeq a$, 那么$a\approx b$.
\end{enumerate}

\end{document}
